

% The results indicate that with increasing proportion of late successional forest in the surrounding landscape, the seed richness increases (fig. \ref{fig:stats} a)), and that with increasing proportion of forest cover in the surrounding landscape, the proportion of generalist seeds in the seed rain decreases (fig. \ref{fig:stats} b)). Also, we found a significantly higher seed richness and proportion of biotically dispersed seeds in wet tropical forest compared to dry tropical. Additionally, we found a significantly higher proportion of abiotically dispersed seeds in the dry tropical forest compared to wet tropical forest (fig. \ref{fig:stats} c)).

In conclusion, this study demonstrates that the natural regeneration of tropical secondary forests is influenced by the surrounding forest landscape. Specifically, a higher proportion of late successional forest in the landscape increases woody plant seed richness, highlighting the role of mature forests as reservoirs of species diversity. In contrast, forest cover and connectivity did not show a significant relationship with seed richness, nor did forest cover, connectivity, or successional stage significantly affect dispersal modes.

We found that increased forest cover was associated with a decrease in generalist species, which may create favorable conditions for shade-tolerant species. However, no significant relationships were detected between forest connectivity or successional stage and generalist species, and neither pioneer nor shade-tolerant species showed significant responses to forest cover, connectivity, or successional stage.

Forest type emerged as a key factor shaping both seed richness and dispersal mode. Wet tropical forests showed higher seed richness and a greater proportion of biotically dispersed seeds, whereas dry tropical forests were dominated by abiotically dispersed seeds. Functional guilds did not differ significantly between the two forest types.

These findings highlight the importance of considering both successional stage of surrounding forests and the forest type when designing conservation and restoration strategies. Placing regenerating forests near late successional areas and tailoring conservation approaches to local ecological conditions can significantly improve the success of natural forest recovery.

